% Consolidated proofs from the 26 October continuation; citations normalized.
\section{Теорема о минимальной степени среза}\label{sec:jets}
Положим $V=\gamma(u)-t^\mu X$. Для любого $P\circ\gamma=u$ имеется полиномиальный срез
\[
S_{P,\mu}=\frac{u-P(V)}{t^\mu},\qquad D_\mu(S_{P,\mu})=1.
\]
Полиномиальность следует из точной делимости числителя; знаменатель не означает рациональную функцию. Эта классическая конструкция восходит к van den Essen--van Rossum \cite{vdER}. Наше дальнейшее вычисление уточняет задачу минимизации степени.

\subsection*{Теорема 3.1. Полнота сведения}
Минимум степени среди всех полиномиальных срезов равен минимуму $\deg S_{P,\mu}$ по всем полиномам параметра $P$.

\textit{Доказательство.} Введём веса
$\wt(t)=1$, $\wt(u)=0$, $\wt(X_i)=-\mu$. Дифференцирование повышает вес на $\mu$. Если $D_\mu(s)=1$, компонента $s_{-\mu}$ веса $-\mu$ тоже является срезом; её обычная степень не больше степени $s$.
После обращения $t$:
\[
A_t=k[t,t^{-1},u,V],\qquad D_\mu=t^\mu\partial_u.
\]
Элемент $t^\mu s_{-\mu}-u$ инвариантен и имеет вес нуль. Поскольку все $V_i$ имеют вес нуль, он равен $-P(V)$ для единственного $P\in k[V]$. Редукция точного полиномиального равенства по модулю $t$ даёт $P\circ\gamma=u$.\hfill$\square$

\subsection*{Теорема 3.2. Точная формула и конечный критерий}
С соглашением $\deg0=-\infty$ имеем
\[
\boxed{\deg S_{P,\mu}=\max_{\alpha\ne0}
\bigl\{\deg_u((\partial^\alpha P)\circ\gamma)
+(\mu+1)|\alpha|-\mu\bigr\}.}
\]
Действительно, формула Тейлора даёт
\[
S_{P,\mu}=\sum_{\alpha\ne0}\frac{(-1)^{|\alpha|+1}}{\alpha!}
t^{\mu(|\alpha|-1)}X^\alpha(\partial^\alpha P)(\gamma(u)).
\]
Различным $\alpha$ соответствуют различные мономы по $X$, поэтому сокращения между ними невозможны. Производная порядка $\deg P=r$ даёт нижнюю оценку $(\mu+1)r-\mu$.

Следовательно, срез степени $\le d$ существует тогда и только тогда, когда разрешима конечная линейная система
\[
\begin{gathered}
\deg P\le\floor{\frac{d+\mu}{\mu+1}},\qquad P\circ\gamma=u,\\
[u^j](\partial^\alpha P)(\gamma(u))=0\quad
\text{при }j>d-(\mu+1)|\alpha|+\mu,\quad\alpha\ne0.
\end{gathered}
\]
Она линейна по коэффициентам $P$. Для $n=3$ число неизвестных не превышает
$\binom{\lfloor(d+\mu)/(\mu+1)\rfloor+3}{3}$.
Это полный критерий, а не ограничение поиска выбранным видом среза.

\subsection*{Теорема 3.3. Все достаточно большие кратности}
Пусть
\[
M=\max_i\deg f_i,\qquad
r_\gamma=\min\{\deg P:P\circ\gamma=u\},\qquad
m_\gamma(\mu)=\min_{D_\mu(s)=1}\deg s.
\]
Тогда
\[
\boxed{m_\gamma(\mu)=(\mu+1)r_\gamma-\mu\quad\text{при }\mu\ge M-1.}
\]
Для непостоянного $P$ степени $r$ производная порядка $q$ после подстановки имеет степень не больше $M(r-q)$. При $\mu+1\ge M$:
\[
M(r-q)+(\mu+1)q-\mu\le(\mu+1)r-\mu.
\]
Ненулевая производная порядка $r$ достигает этой границы. Минимизация по $P$ завершает доказательство. Кроме того, $m_\gamma(\mu)$ не убывает: каждое выражение в максимуме имеет неотрицательный наклон $|\alpha|-1$.

Эта метрика зависит от исходных координат. Она не является инвариантом сопряжения LND и не может без дополнительного аргумента служить препятствием полиномиальности ядра.

\section{Точные вычисления для гармонических моделей}\label{sec:computed}
Для $\mu=1$ получены следующие результаты. Нижние границы удостоверены рациональными дуальными сертификатами, верхние --- явными срезами.

\begin{center}\small\renewcommand{\arraystretch}{1.22}
\begin{tabular}{@{}lrrr@{}}\toprule
Модель & $m_\gamma(1)$ & Неизвестных & Ненулевых весов\\\midrule
$H(3,4,5)$ &7&20&11\\
$H(3,5,7)$ &11&56&35\\
$H(4,5,7)$ &12&84&41\\
$H(3,7,8)$ &13&84&41\\
$H(3,10,11)$ &19&220&111\\
$H(3,13,14)$ &25&455&225\\
$H(5,7,8)$ &11&56&29\\
$H(7,10,11)$ &15&120&61\\\bottomrule
\end{tabular}\end{center}
Последние два столбца относятся к системе, исключающей степень на единицу меньше минимума. Первые две строки воспроизводят прежние результаты более компактными сертификатами.

Если система записана $A c=b$, нижний сертификат есть рациональный вектор $\lambda$ с $\lambda^{\mathsf T}A=0$ и $\lambda^{\mathsf T}b=1$. Это исключает решение над любым полем характеристики нуль, а не только над $\Q$. Независимый проверяющий скрипт заново строит столбцы и выполняет вычисления с \file{Fraction}; символьный решатель для проверки не нужен.

\subsection*{$H(3,7,8)$: меньшая степень ретракции не всегда лучше}
Тождества Чебышёва дают
\[
C_7C_8=C_{15}+C_1,\qquad C_5(C_3)=C_{15},\qquad
P_5=yz-C_5(x)=yz-x^5+5x^3-5x.
\]
Полученный срез степени $13$ оптимален при $\mu=1$. Однако минимальная степень самого полинома параметра равна $4$; можно взять
\[
\begin{split}
P_4={}&x^3z+x^3-x^2yz-xy^2-2xz-x\\
&+y^3+2yz-2y.
\end{split}
\]
При $\mu=1$ этот $P_4$ даёт срез степени $19$. Таким образом, минимизировать $\deg P$ и минимизировать $\deg S_P$ --- разные задачи.

Полная зависимость имеет вид
\[
\boxed{\begin{aligned}
m_{378}(1)&=m_{378}(2)=13,\quad m_{378}(3)=17,\quad m_{378}(4)=21,\\
m_{378}(5)&=22,\quad m_{378}(6)=23,\quad
m_{378}(\mu)=3\mu+4\quad(\mu\ge7).
\end{aligned}}
\]
Для $\mu\le4$ достаточно $P_5$, начиная с $\mu=5$ --- $P_4$. Нижние границы при $1\le\mu\le7$ проверены точно; последний тест одновременно исключает $\deg P\le3$. Далее применяется теорема 3.3.

\subsection*{$H(4,5,7)$: формула для всех кратностей}
Здесь
\[
P_3=(x^2-2)z-y^3+3y,\qquad
\boxed{m_{457}(\mu)=\max\{12,\mu+9,2\mu+3\}.}
\]
Первые, вторые и третьи производные $P_3$ вдоль кривой имеют максимальные степени $11$, $7$, $0$. Это даёт верхнюю оценку. При $\mu=1,2,3$ нижняя оценка $12$ следует из сертификата при $\mu=1$ и монотонности.

Ограничение полиномов степени $\le3$ на кривую инъективно: соответствующий минор размера $20\times20$ имеет определитель $267386$. Поэтому $P_3$ единственен в этой степени. При $\mu=4,5$ любой $P$ степени $\ge4$ даёт степень среза хотя бы $3\mu+4>\mu+9$. При $\mu\ge6$ действует общая теорема с $M=7$, $r_\gamma=3$.

\subsection*{Бесконечное семейство верхних границ}
Для нечётных $a,m\ge3$ положим $b=(ma-1)/2$, $c=(ma+1)/2$. Индексы попарно взаимно просты, и $P=yz-C_m(x)$ восстанавливает параметр. Точная степень его среза:
\[
\deg S_{P,\mu}=\max\{a(m-1)+1,(\mu+1)m-\mu\}.
\]
Для $(a,b,c)=(3,3n+1,3n+2)$ получаем $m_\gamma(1)\le6n+1$. Равенство удостоверено при $n=1,2,3,4$; общая нижняя граница пока не доказана.

\section{Ядро и формальная тривиальность}\label{sec:formal}
Пусть $I=\ker(k[Y_1,\ldots,Y_n]\to k[u])$ --- идеал вложения, $B_\mu=\ker D_\mu$. Через разложение Тейлора получается классическое по своей природе описание
\[
\boxed{B_\mu=k[t,V,I(V)/t^\mu]\simeq k[Y][t,I/t^\mu].}
\]
Для $\mu=1$ соответствующая конструкция описана в \cite[Proposition 5.4, Corollary 5.5]{vdER}. Она не утверждает полиномиальности $B_\mu$.

Элементарный ключ: если все производные $F$ порядков $<q$ исчезают вдоль $\gamma$, то $F\in I^q$. Для доказательства выбираем $P\circ\gamma=u$ и $g_i=Y_i-f_i(P(Y))$. Разложение Тейлора вокруг $\gamma(P(Y))$ показывает сначала $I=(g_1,\ldots,g_n)$, затем требуемую принадлежность. Обратное включение следует из правила Лейбница. Весовая компонента $t^jF(V)$ полиномиальна при $j<0$ ровно тогда, когда $F\in I^{\lceil-j/\mu\rceil}$; это даёт формулу ядра.

\subsection*{Теорема 5.1. Все конечные утолщения тривиальны}
Для любого вложения, любого $\mu\ge1$ и $N\ge1$:
\[
\boxed{B_\mu/(t^N)\simeq(k[t]/(t^N))^{[n]}.}
\]
Изоморфизмы можно выбрать совместимыми по $N$, поэтому
\[
\boxed{\widehat B_\mu^{\,(t)}\simeq k[p,q_1,\ldots,q_{n-1}][[t]].}
\]
Справа стоят ряды по $t$ с полиномиальными коэффициентами. Степени коэффициентов по $p,q$ могут расти. Это не утверждение об обычном полиномиальном кольце $k[[t]][p,q]$.

\textit{Доказательство.} Положим $a(u)=\nabla P(\gamma(u))$, $v(u)=\gamma'(u)$. По правилу цепочки $a v=1$. Над PID $k[u]$ модуль строк, аннулирующих $v$, свободен ранга $n-1$; выберем базис $\ell_1,\ldots,\ell_{n-1}$. Вместе с $a$ эти строки образуют обратимую над $k[u]$ матрицу.

Строку $\ell$ можно явно поднять в идеал кривой:
\[
F_\ell(Y)=\sum_i\ell_i(P(Y))\bigl(Y_i-f_i(P(Y))\bigr).
\]
Имеем $F_\ell\in I$ и $\nabla F_\ell(\gamma)=\ell-(\ell v)a=\ell$. Определим
\[
p=P(V),\qquad q_j=-F_{\ell_j}(V)/t^\mu,\qquad s=(u-P(V))/t^\mu.
\]
Это полиномы; $p,q_j$ инвариантны, $D_\mu(s)=1$. При $t=0$ замена $(u,X)\mapsto(p,q,s)$ равна
\[
(u,X)\longmapsto(u,\ell_1(u)X,\ldots,\ell_{n-1}(u)X,a(u)X)
\]
и обратима.

После композиции с её полиномиальной обратной при $t=0$ получаем эндоморфизм $\Phi\equiv\mathrm{id}\pmod{t^\mu}$. На $A/(t^N)$ оператор $E=\mathrm{id}-\Phi$ повышает $t$-адический порядок на $\mu$, поэтому $E^L=0$ при $L=\lceil N/\mu\rceil$. Обратное отображение есть конечная сумма $\sum_{j=0}^{L-1}E^j$. Это аддитивное обратное к гомоморфизму колец, следовательно, оно автоматически мультипликативно.

Итак,
\[
A/(t^N)=(k[t]/(t^N))[p,q_1,\ldots,q_{n-1},s],\qquad\overline D_\mu=\partial_s.
\]
Глобальная теорема о срезе даёт $A=B_\mu[s]$. Поэтому ядро редуцированного дифференцирования равно $B_\mu/(t^N)$. Те же полиномы $p,q$ работают для всех $N$, обеспечивая совместимость.\hfill$\square$

\textbf{Следствие для стратегии.} Абстрактный тип одного конечного утолщения слоя $t=0$ не может отличить эти ядра от полиномиального. Нужно контролировать глобальную алгебраизацию, степени обратного отображения или дополнительные структуры. Конечные сертификаты степенных ограничений по-прежнему полезны: они учитывают исходные координаты, которых нет в абстрактном кольце утолщения.

\section{Ограничение класса координат для кубического семейства}\label{sec:uniform}
Пусть $\gamma_n=(C_3,C_{3n+1},C_{3n+2})$, $n\ge1$. Над $\C$ не существует полиномиальной координаты $H$ с $H\circ\gamma_n=u$ ни одной из форм
\[
\boxed{H=Q(x,y)+M(x)z\quad\text{или}\quad H=Q(x,z)+M(x)y.}
\]
Степени $Q$ и $M$ не ограничиваются. Запрет относится и к абсолютной координате: сохранять $x$ в полном автоморфизме не требуется.

\subsection*{Доказательство относительной части}
Пишем $U_j=U_j(x/2)$ для полиномов Чебышёва второго рода. В свободном модуле с базисом $1,u,u^2-2$ над $\C[x]$, где $u^3-3u=x$, имеем
\[
y=U_nu-U_{n-1}(u^2-2),\qquad
z=-U_{n-1}u+U_n(u^2-2).
\]
Для первой формы редуцируем $Q$ по моническому соотношению
$y^3-3y-C_{3n+1}(x)=0$ до $Ay^2+By+C$. Сравнение коэффициентов даёт
\[
U_{3n}A+U_{2n}M=U_{n-1}.\tag{6.1}
\]
При $M=0$ равенство невозможно по степеням. Если $(x,H,W)$ автоморфизм и $M(\alpha)=0$, полином $Q(\alpha,y)$ должен быть аффинным по $y$. Поэтому $A(\alpha)=0$ и $U_{n-1}(\alpha)=0$.

Все корни $M$ лежат среди вещественных корней $U_{n-1}$. Два наибольших корня $U_{3n}$,
$x_j=2\cos(j\pi/(3n+1))$, $j=1,2$, лежат правее всех этих корней. При $n=1$ корней $U_0$ вообще нет. Подстановка в (6.1) и синусная формула дают
$M(x_1)=1$, $M(x_2)=-1$.
Но для $M=c\prod(x-\alpha)^{e_\alpha}$ отношение $M(x_1)/M(x_2)$ положительно, даже при комплексном $c$. Противоречие.

Для второй формы аналогично получается
\[
U_{3n+1}A+U_{2n}M=U_n.
\]
Два наибольших корня $U_{3n+1}$ лежат правее корней $U_n$ и опять принуждают противоположные значения $M$. Это завершает относительное доказательство.

\subsection*{Переход к абсолютной координате}
При $M\ne0$ полином $H$ является координатой над $\C(x)$. Если он был бы абсолютной координатой, Proposition 3.1 работы Edo--van den Essen--Maubach \cite{EVM} делала бы его координатой над $\C[x]$. Противоречие уже доказанному.

Для $n=1$ запрет совместим с имеющимися в проекте комплексными конструкциями для $H(3,4,5)$: параметрическая компонента должна нелинейно зависеть от обеих поперечных переменных. Это ограничение формы выпрямителя, а не доказательство его отсутствия.

\subsection*{Естественная модификация тоже не проходит}
Рассмотрим
\[
\beta_n=(x,C_n(x)z-y^2+2,y),\qquad
\beta_n\circ\gamma_n=(C_3,C_2,C_{3n+1}).
\]
На каждом корне $\alpha$ многочлена $C_n$ центр $H(\alpha,2-y^2,y)$ не может быть аффинным по $y$, если $H$ восстанавливает параметр целевой тройки. Иначе в факторе по $u^3-3u-\alpha$ получаем
\[
u=\lambda\bigl(U_n(\alpha/2)u-U_{n-1}(\alpha/2)(u^2-2)\bigr)+\nu.
\]
Здесь $U_{n-1}(\alpha/2)\ne0$, поскольку
$C_n(x)^2+(4-x^2)U_{n-1}(x/2)^2=4$. Сравнение коэффициентов сначала даёт $\lambda=0$, затем противоречие. Это доказательство не имеет ограничения степени $H$.

При нечётном $n$ целевая тройка не является вложением: $u=\pm\sqrt3$ имеют одинаковые образы. При чётном $n=2m$ цель имеет явную относительную координату. Если
$C_{6m+1}(u)=uP_m(u^2-2)$, $\varepsilon=P_m(1)=(-1)^m$, то
\[
A_m(Y)=\frac{1-\varepsilon P_m(Y)}{Y-1},\qquad
H_0=xA_m(Y)+\varepsilon Z
\]
восстанавливает параметр. Таким образом, даже выпрямляемость цели не исправляет плохой центр данного переноса.

\subsection*{Дополнительный конечный результат для $H(3,7,8)$}
Исключены координаты $Q(x,L)+M(x)N$, где $(L,N)$ --- любая постоянная невырожденная линейная замена $(y,z)$, при $\deg M\le10$ и произвольной степени $Q$. После нормировки $L=y+az$ задача сводится к 56 точным условиям на кратности корней $M$. Во всех случаях НОД необходимых полиномиальных условий равен $1$; повторные корни при $a=\pm i/2$ проверены отдельно. Это \textbf{ограниченный результат}: степень $M>10$ здесь не исключена.

\section{H(4,5,7): барьер степени четыре в относительном классе}\label{sec:h457}
Пусть $F\circ(C_4,C_5,C_7)=u$ и $\deg F\le4$. Тогда:
\begin{itemize}
\item единственный такой $F$, являющийся координатой $k(x)[y,z]$, есть $F_0=(x^2-2)z-y^3+3y$;
\item ни один такой $F$ не является координатой $k(y)[x,z]$ или $k(z)[x,y]$;
\item $F_0$ не является координатой над $k[x]$: его производные по $y,z$ одновременно исчезают при $x^2=2$, $y^2=1$.
\end{itemize}
Следовательно, для любого метода, дополняющего одну из пар $(x,F)$, $(y,F)$, $(z,F)$ до автоморфизма, обязательно $\deg F\ge5$.

Доказательство исчерпывает всё пространство полиномов степени $\le4$. Из 35 мономов матрица ограничения на кривую имеет ранг 28. Все полиномы параметра образуют аффинное пространство $F_0+\langle K_0,\ldots,K_6\rangle$. В архиве полностью выписаны семь $K_i$; полнота проверена ненулевыми минорами.

Далее применяется теорема Юнга--ван дер Кулка: координата плоскости степени 3 имеет вид $P_3(L)+cM$, а степени 4 --- либо $P_4(L)+cM$, либо $A(M+Q_2(L))^2+bL+c$ с соответствующими ненулевыми коэффициентами. Последовательное сравнение коэффициентов во всех трёх проекциях исключает варианты. В последней ветви вычисляется ненулевой результант $6561301239$; имеется явное тождество Безу с постоянной $729033471$.

Верификатор проверяет 35 точных условий, а отдельный аудит повторяет ранги, граничные случаи и устранение параметров. Структурный источник для нормальных форм --- классическая теорема и её лемма о делимости степеней; удобный первичный текст указан в \cite{Chau}.

\textbf{Граница утверждения.} Оно не исключает абсолютную координату степени 4, завершение которой меняет все три исходные координаты, и не исключает относительные координаты степени $\ge5$.

\section{Полная фильтрация одношагового переноса}\label{sec:transfer}
Пусть $R=k[x]$, $A=R[y,z]$, $(H,W)\in\Aut_R R[Y,Z]$ и
\[
d=H_YW_Z-H_ZW_Y\in k^*,\qquad
\beta=(p(y)+mz,y),\quad p\in R[y],\ 0\ne m\in R.
\]
Положим $h=H\circ\beta$, $w=W\circ\beta$, $D=h_z\partial_y-h_y\partial_z$.
Вынесем содержание по непостоянным мономам:
\[
c=h(x,0,0),\qquad h=c+g h_0,
\]
где $g\in R$ --- монический НОД коэффициентов всех непостоянных по $y,z$ мономов $h$.

\subsection*{Теорема 8.1}
Имеем $g\mid m$, $B=\ker D=k[x,h_0]$. Для каждого неприводимого $\pi\mid m_1=m/g$ определим
\[
E_\pi=v_\pi(m_1),\qquad e_\pi=\deg_y(h_0\bmod\pi)\ge1.
\]
Редукция $h_0$ здесь не зависит от $z$. Для $\mathcal F_n=\ker D^{n+1}$:
\[
\boxed{D^n(\mathcal F_n)=
\left(g^n\prod_{\pi\mid m_1}\pi^{E_\pi(n-\lfloor n/e_\pi\rfloor)}\right)B.}
\]
В частности,
\[
\boxed{\pl(D)=D(A)\cap B=
\left(g\prod_{e_\pi\ge2}\pi^{E_\pi}\right)B.}
\]
Все кратности исключительного делителя сохраняются. Хорошие компоненты $e_\pi=1$ исчезают из plinth-идеала примитивного дифференцирования; множитель $g$ учитывается отдельно.

\subsection*{Основные шаги доказательства}
\textit{1. Ядро и нормировка.} Над $k(x)$ пара $(h,w)$ является системой координат, $D(w)=dm$. Отсюда $g\mid m$ и общее ядро равно $k(x)[h_0]$. Редукция $h_0$ по любому простому $\pi\in R$ непостоянна; поэтому $R[h_0]$ насыщено по $\pi$ в $A$. Последовательное сокращение знаменателей даёт $A\cap k(x)[h_0]=R[h_0]$.

\textit{2. Точное уравнение модификации.} Пусть $(\mathcal Y,\mathcal Z)$ --- обратное к $(H,W)$ отображение. Полином
\[
N(T,V)=\frac{\mathcal Y(c+gT,V)-p(\mathcal Z(c+gT,V))}{g}
\]
действительно лежит в $R[T,V]$, причём $\gcd(N,m_1)=1$; это доказывается редукцией по $g$ и обратимостью центра над конечным факторкольцом. Затем
\[
A\simeq R[T,V,z]/(m_1z-N(T,V)),\qquad T=h_0,\quad V=w.
\]
Полиномиальность $N$ существенна: формальное деление без её проверки оставило бы пробел.

\textit{3. Моническая модель над каждым простым.} Локализуем только $R$: $O=R_{(\pi)}$, коэффициентное кольцо $O[T]$. Над $O/(\pi^{E_\pi})$ центр имеет вид $(G(y),J(y))$, порождающий всё кольцо полиномов по $y$. Оно свободно ранга $e_\pi$ над кольцом по $G$; доказательство получается повторением обычного деления и конечным подъёмом по степеням $\pi$.
Характеристический полином умножения на $J$ даёт монический $P_\pi(T,V)$ степени $e_\pi$ по $V$. Уравнения центра и $P_\pi$ порождают один идеал по модулю $\pi^{E_\pi}$, что приводит к точной модели
\[
P_\pi(h_0,w)=\pi^{E_\pi}v.
\]
Здесь сохраняются условия, зависящие от $T$: замена $O[T]$ полем рациональных функций по $T$ была бы недостаточна.

\textit{4. Базис и старший коэффициент.} Мономы $w^iv^j$, $0\le i<e_\pi$, образуют базис над $O[h_0]$. Их $D$-степени равны $i+e_\pi j$ и попарно различны. Поэтому старшие члены не сокращаются. При $n=i+e_\pi j$ получаем показатель $E_\pi(n-j)$. Пересечение локальных условий по всем простым $R$, с использованием китайской теоремы об остатках для построения предобразов, даёт глобальную формулу.

Полное алгебраическое доказательство, включая делимость, монический подъём и склеивание, приведено в \file{relative_lnd/PROOF.md}; отдельный аудит находится в \file{literature/independent_image_ideal_review.md}. Дополнительно проверены 16 примеров и 28 независимых точных условий.

\subsection*{Пример с кратным делителем и нелинейным подъёмом}
Пусть
\[
h=y^2+xy^3+x^2z,\qquad D=x^2\partial_y-(2y+3xy^2)\partial_z.
\]
Здесь $E=2$, $e=2$. Для $v=y^4+xyz-z$:
\[
Dv=2y+xh,\qquad D^2v=2x^2,\qquad
z=-(1+xy)v+hy^2.
\]
Последовательность идеалов образов имеет показатели $0,2,2,4,4,6,6,\ldots$. Один найденный элемент образа даёт лишь включение; равенство всего идеала следует из теоремы.

\subsection*{Следствия для чистой теории LND}
Для $D_{M,e}=x^M\partial_y-e y^{e-1}\partial_z$, $M\ge1$, $e\ge2$:
\[
I_n(D_{M,e})=(x^{M(n-\lfloor n/e\rfloor)}).
\]
Показатель $M$ восстанавливается из кратности простого делителя plinth-идеала, а $e$ --- как первое $n$ с $I_n\supsetneq I_1^n$. Поэтому разные пары $(M,e)$ дают несопряжённые LND. Само семейство классическое; приоритет этого способа записи следствия не заявляется.

Для $h=(x^2-2)z-y^3+3y$, возникающего у $H(4,5,7)$:
\[
\ker D=k[x,h],\qquad I_n(D)=(x^2-2)^{n-\lfloor n/3\rfloor}k[x,h].
\]
Это точная информация о данном LND, но не препятствие для всех возможных выпрямителей кривой.

\section{Последовательные модификации и группы Кремоны}\label{sec:iterated}
Определим
\[
r_{-1}=z,\quad r_0=y,\quad r_j=xr_{j-2}-r_{j-1}^2,\qquad j\ge1.
\]
Для фиксированного $\ell\ge1$ положим $p=r_\ell$, $w=r_{\ell-1}$ и
\[
D_\ell=(-1)^\ell\bigl(p_y\partial_z-p_z\partial_y\bigr).
\]
Тогда $D_\ell(w)=x^\ell$, $\ker D_\ell=B=k[x,p]$. Дифференцирование неприводимо: общий делитель компонент должен делить $x^\ell$, а $p\bmod x=-y^{2^\ell}$ имеет ненулевую производную по $y$.

На открытом множестве $x\ne0$ каждая операция обратима; знаменатели появляются при продолжении через $x=0$. Это конкретная связь с преобразованиями Кремоны: биррациональная обратимость проста, а полиномиальность и длина разложения контролируются исключительным делителем.

\subsection*{Теорема 9.1. Вся фильтрация}
Для $I_n(D)=D^n(\ker D^{n+1})$:
\[
\boxed{I_n(D_\ell)=x^{a_\ell(n)}B,\qquad
a_\ell(n)=\ell n-\sum_{j=1}^{\ell}\floor{\frac n{2^j}}.}
\]
Доказательство даёт явный базис. Положим $v_j=r_{\ell-1-j}$ для $0\le j\le\ell$, $v_{-1}=p$. Тогда
\[
v_j^2=xv_{j+1}-v_{j-1}\quad(0\le j<\ell),\qquad \deg_{D_\ell}v_j=2^j.
\]
По моническим соотношениям нормальные мономы имеют вид
\[
s_n=\prod_{j=0}^{\ell-1}v_j^{\varepsilon_j}\,v_\ell^q,
\qquad n=\sum_{j=0}^{\ell-1}\varepsilon_j2^j+q2^\ell,
\quad\varepsilon_j\in\{0,1\}.
\]
Они образуют $B$-базис и имеют попарно разные $D$-степени. В выражении через $w$ старший коэффициент $v_j$ равен $x^{-(2^j-1)}$. Применение $D_\ell^n$ и двоичная запись $n$ дают указанную сумму целых частей; множитель $n!$ является единицей в $k$.

При $\ell=3$ показатели для $0\le n\le8$:
\[
0,\ 3,\ 5,\ 8,\ 9,\ 12,\ 14,\ 17,\ 17.
\]
При $\ell\ge2$ дифференцирование $D_\ell$ не сопряжено ни с одним неприводимым одношаговым переносом теоремы 8.1, даже при изменении базовой координаты. Действительно, единственный простой делитель plinth-идеала и его кратность $\ell$ сохраняются при любом сопряжении. Для одношагового класса показатель $I_2$ на этом простом был бы $\ell$ при $e=2$ или $2\ell$ при $e>2$, тогда как здесь он равен $2\ell-1$.

В частности, $D_\ell$ при $\ell\ge2$ \textbf{не триангулируемо}. Если бы в некоторых координатах оно имело вид $a(X)\partial_Y+b(X,Y)\partial_Z$ с $a\ne0$, оно принадлежало бы одношаговому классу с $h=a(X)Z-\int b(X,Y)\,dY$. Если первая координата имеет ненулевую постоянную производную, plinth равен всему ядру; случай $a=0$ при неприводимости также даёт постоянный перенос. Оба случая несовместимы с нашим plinth. Само существование таких нетриангулируемых примеров не объявляется новым; здесь дан явный способ распознать их двумя идеалами.

\subsection*{Ассоциированное градуированное кольцо}
Те же монические соотношения дают более полное описание:
\[
\boxed{\gr_{D_\ell}A\simeq
B[\xi_0,\ldots,\xi_\ell]/(\xi_j^2-x\xi_{j+1}:j<\ell),\quad
\deg\xi_j=2^j.}
\]
После редукции по $x$ нильрадикал $N=(\xi_0,\ldots,\xi_{\ell-1})$ удовлетворяет $N^\ell\ne0$, $N^{\ell+1}=0$, а $N/N^2$ свободен ранга $\ell$ над $k[p,\xi_\ell]$.

Показательный пример: $D_3$ и $xD_2$ имеют одинаковый plinth $(x^3)$, и показатели $I_n$ совпадают при $n\le7$. При $n=8$ они равны соответственно $17$ и $18$. Полная фильтрация видит различие, которое первые семь идеалов не обнаруживают.

Это иллюстрация дополнительной информации фильтрации, а не новая трудная
задача сопряжения: сами поля уже различаются неприводимостью.
При $\ell=2$ нетриангулируемость также известна из
\cite[Example 4.19]{DG}; новизна общего семейства и формул здесь
не устанавливается.

\subsection*{Связь с актуальным вопросом о сложности}
В журнальной статье Dasgupta--Gaifullin \cite[Theorem 4.12]{DG} universal level связан со сложностью биррационального представления формулой $\operatorname{ulev}(D)=c(\alpha)+1$. Их Question 4.21 относится к кандидату, построенному тремя чередующимися модификациями и перестановками.

Наша рекурсия позволяет исследовать соответствующий тип конструкции через точную фильтрацию. \textbf{Но равенство $c(\varphi)=3$ здесь не доказано}: для него требуется исключить все более короткие допустимые разложения, а не только канонический одношаговый класс. Возможность извлечь такую нижнюю оценку из $\gr_DA$ --- содержательная следующая задача.

\section{Почему относительный ML не различает наши кривые}\label{sec:relativeML}
Пусть $A=k[x_1,\ldots,x_n]$, $I\subsetneq A$ и $\dim(A/I)\le n-2$. Для любого $r\ge1$:
\[
\boxed{\bigcap_{\substack{0\ne D\in\operatorname{LND}(A)\\D(A)\subset I^r}}\ker D=k.}
\]
Более того, алгебра, порождённая объединением этих ядер, равна $A$.

Для каждого $i$ существует ненулевой $f_i\in I\cap k[x_1,\ldots,\widehat{x_i},\ldots,x_n]$: иначе $n-1$ образов координат были бы алгебраически независимы в кольце размерности $\le n-2$. Возьмём
\[
D_i=f_i^r\partial_{x_i},\qquad
\ker D_i=k[x_1,\ldots,\widehat{x_i},\ldots,x_n].
\]
Эти LND удовлетворяют $D_i(A)\subset I^r$, пересечение их ядер равно $k$, а объединение содержит все координаты. Это доказывает оба утверждения.

Следовательно, кандидат с более слабым условием $D(I)\subset I$ тоже всегда даёт $k$ для кривой в $\A^3$. Даже требование поточечно фиксировать произвольное конечное утолщение не помогает.

Для гармонических моделей можно взять, например,
\[
f_x=C_c(y)-C_b(z),\quad
f_y=C_c(x)-C_a(z),\quad
f_z=C_b(x)-C_a(y).
\]
Механизм реплик согласуется с более сильными результатами о гибкости дополнений, см. Flenner--Kaliman--Zaidenberg \cite{FKZ}. Поэтому здесь исправляется исследовательская стратегия, а не объявляется новый нетривиальный инвариант.

